%==============================================================================
% ABSTRACT
%==============================================================================
\chapter*{Abstract}
\addcontentsline{toc}{chapter}{Abstract}
\markboth{Abstract}{Abstract}

\noindent
\textbf{\thesismajor}\\[4pt]
\textbf{Student:} \thesisauthor \qquad \textbf{Adviser:} \thesisadviser

\vspace{1em}

Low-Dose Computed Tomography (LDCT) has become a cornerstone of modern clinical diagnostics,
enabling substantial reductions in patient radiation exposure without sacrificing diagnostic
utility. However, the inherent trade-off between dose reduction and image quality
introduces complex, non-stationary noise degradations and streak artefacts that critically
impair radiological assessment. Accurate, automated evaluation of perceptual image quality
is therefore essential for quality assurance in high-throughput clinical pipelines, yet
the field lacks reliable, no-reference methods specifically designed for medical CT imaging.

This thesis proposes \textbf{CT-MUSIQ}, a novel no-reference image quality assessment
(NR-IQA) framework that adapts the Multi-Scale Image Quality Transformer (MUSIQ)
architecture to the specific biophysical and perceptual constraints of brain CT imaging.
CT-MUSIQ accepts a single axial CT slice as input, applies the clinical brain soft-tissue
windowing transformation (level~40~HU, width~80~HU), and produces a continuous perceptual
quality score in the range~$[0,4]$ that closely aligns with consensus radiologist Likert
scale annotations.

The proposed architecture introduces four key innovations over the vanilla MUSIQ baseline:
(1)~a \textbf{multi-scale spatial pyramid} operating at two resolutions ($224 \times 224$
and $384 \times 384$ pixels), producing 193 patch tokens per image that simultaneously
capture global structural context and fine-grained noise texture; (2)~a
\textbf{hash-based positional encoding} scheme that injects scale, row, and column
identity into each token without the sequence-length restrictions of standard sinusoidal
encodings; (3)~a \textbf{scale-consistency Kullback--Leibler divergence regularisation}
loss that penalises disagreement between per-scale quality predictions and the global
head, preventing the model from learning scale-specific shortcuts on the small training
dataset; and (4)~a \textbf{pairwise margin ranking loss} component that explicitly
optimises the rank-based correlation metrics used in the LDCTIQAC 2023 grand challenge
evaluation protocol.

The system is trained on the LDCTIQAC 2023 dataset comprising 1,000 brain CT slices
annotated by expert radiologists, using a two-stage fine-tuning strategy with ViT-B/32
pretraining, mixed-precision training, exponential moving average weight averaging, and
gradient clipping, designed to fit within 6~GB of consumer GPU memory.

Experimental evaluation on the held-out test set of 100 images demonstrates that CT-MUSIQ
achieves a Pearson Linear Correlation Coefficient (PLCC) of \textbf{0.9498}, a Spearman
Rank-Order Correlation Coefficient (SROCC) of \textbf{0.9488}, a Kendall Rank-Order
Correlation Coefficient (KROCC) of \textbf{0.8249}, and a composite Aggregate score of
\textbf{2.7235}, placing it second among all methods evaluated relative to the published
LDCTIQAC 2023 leaderboard. An ensemble baseline (Agaldran Combo: Swin-T + ResNet-50)
achieves an Aggregate of 2.7682, outperforming the published leaderboard top submission
of 2.7427, with CT-MUSIQ remaining highly competitive. Ablation studies confirm the
individual contributions of each architectural component, with the scale-consistency KL
loss providing the most significant regularisation benefit on the limited training data.

These results demonstrate that a carefully adapted Vision Transformer architecture with
domain-specific windowing, multi-scale tokenisation, and scale-consistency regularisation
can achieve state-of-the-art automated perceptual quality assessment for brain CT imaging,
providing a computationally efficient and clinically interpretable foundation for
integration into real-world Picture Archiving and Communication Systems (PACS) pipelines.

\vspace{0.5cm}
\noindent\textbf{Keywords:} Computed Tomography, Low-Dose CT, No-Reference Image Quality
Assessment, Vision Transformer, MUSIQ, Perceptual Quality, Brain CT, Multi-Scale Learning,
KL Divergence Regularisation.
